\cgeentry{cge_equation}{User equations: new relations among the model's variables}

\cgesec{Description}

\cmd{cge_equation} adds an equation to the calibrated model, written in a small
language of the model's variables and parameters, for one element or for every
element of a set. With \manref{cge_variable}, which adds variables, it extends
the model without changing its code: Stata compiles the equations for its
solver, and \manref{cge_gams} writes them into the GAMS program, so that GAMS
solves the same model. \cmd{cge_equation, list} lists the user variables and
equations, and \cmd{cge_equation, clear} removes them.

\cgesec{Quick start}

Government purchases of services at \cmd{LAM} times their base level, with
\cmd{LAM} a user variable
\begin{stcode}
cge_variable LAM, value(1)
cge_equation lam: QG(com_ser) = LAM * base(QG(com_ser))
\end{stcode}
Fix \cmd{LAM} and let the scaling factor of government consumption adjust; then
raise it by 10 per cent
\begin{stcode}
cge_closure, swap(LAM GADJ)
cge_shock LAM, value(1.1)
\end{stcode}
An equation for every commodity
\begin{stcode}
cge_variable PIX(c), value(1)
cge_equation pix(c): ln(PIX(c)) = ln(PQ(c)) - ln(CPI)
\end{stcode}
An equation for the commodities with an import tariff only
\begin{stcode}
cge_variable LQM(c)$(tm(c) > 0), value(1)
cge_equation lqm(c)$(tm(c) > 0): QM(c) = LQM(c) * base(QM(c))
\end{stcode}
List the user variables and equations; remove them
\begin{stcode}
cge_equation, list
cge_equation, clear
\end{stcode}

\cgesec{Syntax}

Add an equation
\begin{stsyntax}
\cmdb{cge_equation} \argu{name}[\cmd{(}\argu{i}\cmd{)}[\cmd{$(}\argu{condition}\cmd{)}]]\cmd{:}
\argu{expression} \cmd{=} \argu{expression}
\end{stsyntax}
List the user variables and equations
\begin{stsyntax}
\cmdb{cge_equation}\texttt{,} \opt{}{list}
\end{stsyntax}
Remove the user variables and equations
\begin{stsyntax}
\cmdb{cge_equation}\texttt{,} \opt{}{clear}
\end{stsyntax}
\argu{name} has at most 24 characters, letters, digits and underscores, and
may not be the name of another user equation; a variable defined by
\manref{cge_variable} gives its name to its definition. \argu{i} is an index,
the letter of a set or the name of a user set (\manref{cge_set}), possibly
followed by digits. A condition, \cmd{$(}\argu{condition}\cmd{)}, restricts
the equation to the elements that meet it (from version 0.7.2).

\cgesec{Options}

\begin{optdesc}
\item[\opt{}{list}] lists the user variables, with their definitions or
values, and the user equations. It is the default when only a comma is given.

\item[\opt{}{clear}] removes every user variable and equation. Swaps and
constraints that involve user variables must be removed first, with
\cmd{cge_closure, swap(none)} and \cmd{cge_constraint, clear}.
\end{optdesc}

\cgesec{Remarks and examples}

\cgesubsec{The language}

Expressions are built from
\begin{cgetable}{p{0.34\linewidth}p{0.60\linewidth}}
\toprule
Term & Meaning\\
\midrule\endhead
\bottomrule\endlastfoot
numbers & such as \cmd{2}, \cmd{0.5} or \cmd{1e-3}\\
\cmd{+ - * /} and \cmd{^} & arithmetic; \cmd{^} is the power\\
\cmd{( )} & grouping\\
\cmd{exp()}, \cmd{ln()} or \cmd{log()}, \cmd{sqrt()} & functions\\
\argu{VAR}, \argu{VAR}\cmd{(}\argu{e}\cmd{)}, \argu{VAR}\cmd{(}\argu{e1}\cmd{,} \argu{e2}\cmd{)}
  & a model or user variable at an element, given by its label, such as
    \cmd{com_hlt}, or by the index of the equation or of a sum\\
\argu{par}, \argu{par}\cmd{(}\argu{e}\cmd{)} & a parameter, at its current value\\
\argu{upar}, \argu{upar}\cmd{(}\argu{e}[\cmd{,} \argu{e2}]\cmd{)} & a user parameter (\manref{cge_parameter}), at its current value\\
\cmd{base(}\argu{VAR}[\cmd{(}\argu{e}\cmd{)}]\cmd{)} & the base value of a variable\\
\cmd{sum(}\argu{i}\cmd{,} \argu{expression}\cmd{)} & the sum over the elements of the set of index \argu{i}\\
\cmd{ord(}\argu{i}\cmd{)} & the position of the element of index \argu{i} in its set, from 1\\
\cmd{sum(}\argu{i}\cmd{$(}\argu{condition}\cmd{),} \argu{expression}\cmd{)} & the sum over the elements for which the condition holds\\
\argu{term}\cmd{$(}\argu{condition}\cmd{)} & the term where the condition holds, else 0\\
\end{cgetable}

Variables are written in upper case, with as many elements as they have
indexes, and parameters in lower case. The parameters are \cmd{tm}, \cmd{te},
\cmd{tq}, \cmd{ta}, \cmd{tv}, \cmd{kappa}, \cmd{pwm}, \cmd{pwe}, \cmd{ad},
\cmd{af}, \cmd{qg}, \cmd{qinv}, \cmd{trdg}, \cmd{trdr}, \cmd{trfr},
\cmd{trgr}, \cmd{trrg}, \cmd{tshi}, \cmd{qshi}, \cmd{trshi}, \cmd{trshir},
\cmd{rd}, \cmd{rf}, \cmd{ddebt}, \cmd{fdebt}, \cmd{qlf}, \cmd{wmin},
\cmd{wref}, \cmd{sick} and \cmd{tfy}, as in \manref{cge_shock}, and \cmd{trgg},
transfers within government, and \cmd{hs0}, the base volume of health
services; \cmd{tf}, which has two indexes, and \cmd{mort} cannot be used. A
parameter takes its value in the simulation being solved, with any shock.

The indexes are \cmd{a} activities, \cmd{c} commodities, \cmd{f} factors,
\cmd{h} households, \cmd{d} households and enterprises, \cmd{i} institutions
that receive factor income and \cmd{n} nests of value added, and the names of
user sets (\manref{cge_set}), each possibly followed by digits, as in
\cmd{c2}. An equation
\argu{name}\cmd{(}\argu{i}\cmd{)} is written for every element of the set of
\argu{i}; an index that is neither that of the equation nor that of an
enclosing sum is an error. Sums may contain sums.

\cmd{^} binds most tightly and groups from the right; then come \cmd{*} and
\cmd{/}, and then \cmd{+} and \cmd{-}, which group from the left. A leading
minus sign applies after any power: \cmd{-X^2} is $-(X^2)$.

\cgesubsec{Conditions}

From version 0.7.2, a condition written \cmd{$(}\argu{condition}\cmd{)}, as
in GAMS, restricts a sum, a term, or the domain of an equation or a variable,
\argu{name}\cmd{(}\argu{i}\cmd{)$(}\argu{condition}\cmd{)}. A term is a
number, a reference such as \cmd{PQ(c)}, a function or an expression in
parentheses. Besides the terms of expressions, a condition may use
\begin{cgetable}{p{0.34\linewidth}p{0.60\linewidth}}
\toprule
Term & Meaning\\
\midrule\endhead
\bottomrule\endlastfoot
\cmd{> < >= <= == !=} & comparisons, 1 if true and 0 if not\\
\cmd{&}, \cmd{|}, \cmd{!} & and, or, not; any value other than 0 is true\\
\argu{set}\cmd{(}\argu{e}\cmd{)} & 1 if element \argu{e} belongs to a user set or a model set, as in \cmd{goods(c)}\\
\cmd{sameas(}\argu{i}\cmd{,} \argu{label}\cmd{)} & 1 if the element of index \argu{i} is \argu{label}\\
\cmd{ord(}\argu{i}\cmd{)} & the position of the element in its set\\
\end{cgetable}
Arithmetic is evaluated first, then comparisons, \cmd{!}, \cmd{&} and
\cmd{|}. Comparisons, \cmd{sameas()} and set membership are allowed only in
conditions.

Conditions are evaluated with the variables and parameters at their base
values, when the equation is added and when the model is calibrated again; they
choose the terms and the equation elements, which stay the same in every
scenario. A condition that is missing at the base, such as \cmd{1/tm(c) > 0}
where a tariff is 0, and a domain that no element meets are errors. Stata
reads \cmd{$}\argu{name} on a command line as a global macro, so the
condition is always written in parentheses.

\cgesubsec{Equations and the closure}

Each equation element needs a free variable, so that the closure stays square:
a variable added by \cmd{cge_variable} with \opt{}{value()}, or a variable that
a swap frees. In the example below, the equation determines the new variable
\cmd{LAM}; the swap then fixes \cmd{LAM} and frees the scaling factor of
government consumption, so that \cmd{GADJ} adjusts to satisfy the equation. The
base need not satisfy a new equation: if it does not, \cmd{cge_equation} says
so, and the equation then acts as a shock, moving the solution from the base.

User equations are saved with the model, compiled again when the model is
calibrated again, and solved in every simulation, dynamic run and sensitivity
draw.

\cgesubsec{Example 1: government purchases of services 10 per cent above the base}

The synthetic SAM is mapped and calibrated as in \manref{cge_solve}:
\stlog{eq_lam}
Government purchases of services are 10 per cent above the base, and, with
$\V{GADJ}$ adjusting, so are the other purchases:
\stlog{eq_lam_res}

\cgesubsec{Example 2: indexed equations and sums}

An equation for every commodity defines a price index relative to the CPI in
implicit form, and an equation with sums gives a volume index of household
consumption, the value of consumption over its value at base quantities:
\stlog{eq_index}
The variables could equally be defined by \manref{cge_variable}, with an
expression; the equations show the functions, an index and a sum with
\cmd{base()}.

\cgesubsec{Example 3: listing}

\stlog{eq_list}
\res{r(N_equations)} counts equation elements: one for \cmd{lam}, three for
\cmd{pix} and one for \cmd{cvol}.

\cgesubsec{Example 4: errors, and removing the user equations}

A reference to an element that does not exist is an error, and the user
equations cannot be removed while a swap involves a user variable:
\stlog{eq_err}

\cgesubsec{Example 5: conditions}

Import spending on the commodities that carry a tariff, the tariff revenue of
each, a variable and an equation over the same domain, and government purchases
of goods, with goods a user set:
\stlog{eq_cond}
\cmd{TMV} and \cmd{LQM} exist for the commodities with a tariff only. With the
tariffs removed, the tariff revenue is 0, while the domain stays as it was in the
base. A condition that cannot be evaluated at the base, or that no element
meets, is an error:
\stlog{eq_cond_err}

\cgesec{Stored results}

\cmd{cge_equation} stores the following in \cmd{r()}:
\begin{storedres}
\resgroup{Scalars}
\res{r(N_elements)} & number of elements of the equation, those that meet its condition\\
\res{r(N_variables)} & number of user variables (\opt{}{list})\\
\res{r(N_equations)} & number of equation elements (\opt{}{list})\\
\resgroup{Macros}
\res{r(name)} & name of the equation\\
\end{storedres}

\cgesec{Methods and formulas}

An equation $\ell(x) = r(x)$ with index $i$ is written out for every element
$e$ of the set of $i$, with $i$ replaced by $e$ where it is an argument of a
variable or a parameter; each \cmd{sum(}$j$\cmd{,} $\cdot$\cmd{)} is written
out in the same way over the elements of the set of $j$, inner sums first, so
that
\begin{equation*}
\texttt{sum(}j\texttt{,}\ g\texttt{)} = \sum_{k \in J} g\,[\,j = k\,],
\end{equation*}
where $J$ is the set of index $j$ and $g\,[\,j = k\,]$ is $g$ with $j$ replaced by
$k$. With a condition $q$ (from version 0.7.2),
\begin{equation*}
\texttt{sum(}j\texttt{\$(}q\texttt{),}\ g\texttt{)} = \sum_{k \in J:\ q[\,j = k\,](x^0,\,\theta^0) \neq 0} g\,[\,j = k\,],
\end{equation*}
with $q$ evaluated at the base point $x^0$ and the base parameters $\theta^0$;
a conditional term $g\texttt{\$(}q\texttt{)}$ is $g$ where $q(x^0, \theta^0) \neq 0$
and 0 elsewhere, and an equation with a condition is written out only for the
elements $e$ with $q[\,i = e\,](x^0, \theta^0) \neq 0$. Sums are written out from
the outside in, each element substituted before the sums within, so that an
inner condition may use the index of an enclosing sum.
Each side is then compiled to postfix order, with numbers, variables (as
positions in the stacked vector of variables), parameters, operators and
functions, and \cmd{base()} replaced by the number it denotes. The solver
evaluates both sides at each point, with the parameters of the simulation, and
the equation element is
\begin{equation*}
\ell_e(x) - r_e(x) = 0,
\end{equation*}
a row of the block \cmd{U} of the system. Its derivatives are exact, from
version 0.7.0: the postfix form of each side is differentiated in forward
mode, with respect to the variables that the equation element contains
(section~\ref{sec:jacobian}).
\Chapref{chap:user} describes the compilation and the GAMS program, in which
GAMS differentiates the same equations itself.

\cgesec{Also see}

\begin{alsosee}
\seeentry{cge_variable}{User variables, defined by an expression or determined by user equations}
\seeentry{cge_closure}{Choose the model closure}
\seeentry{cge_shock}{Change a parameter or a fixed variable}
\seeentry{cge_gams}{Write the model as a GAMS program, and solve it in GAMS}
\seechapter{chap:user}{User variables and equations}
\end{alsosee}
